\reportsection{Implications}
\label{sec:implications}

\paragraph{A parameter rule becomes a developer-facing gate.}
The conventional role of a sizing rule is advisory.  Here it is executable.
\code{MeetsBindingTarget} is decidable, and \code{ofField} constructs
parameters that satisfy it.  A concrete configuration can therefore carry the
binding theorem only when its digest width supports the declared security
target and query budget.  The same structure already certifies the salt-space
cardinality used by the hiding development.

\paragraph{The proof stack has clear ownership.}
The binding result can be reviewed one transformation at a time:
\[
\begin{aligned}
\text{inconsistent accepted overlapping openings}
&\Longrightarrow \text{disputed overlap}\\
&\Longrightarrow \text{trace collision}\\
&\Longrightarrow \text{birthday error}\\
&\Longrightarrow \text{concrete parameter rule}.
\end{aligned}
\]
A consuming soundness argument uses binding to treat
accepted views of a committed object as consistent on every shared location.
Where the argument needs to reason about the committed object itself, it may
additionally consume extractability.
Privacy and simulation arguments consume hiding under their stated leakage
model.  The Merkle ROM instance connects the binding property to a collision
event, and the field-specific constructor discharges the numerical
precondition.  This division is useful because later verification can focus on
how these contracts compose, rather than reconstructing them from the hasher
implementation.

\paragraph{The contract should be visible in Rust.}
The remaining implementation proposal is deliberately small.  A shared digest
contract should expose conservative output capacity, and a salt-space contract
should expose cardinality.  Generic code over \code{H: MerkleHasher} can then
check
\[
\code{H::Digest::BIT\_SIZE}\ge\secpar+1+2h
\qquad\text{and}\qquad
\code{H::Salt::CARDINALITY\_BITS}\ge\secpar.
\]
This makes the one-element BabyBear binding and salt choices visible before a
commitment is used.  A hiding-capable implementation must additionally connect
its randomness source and leaf-hash construction to the distribution assumed
by the hiding game.
\status{analytical; proposed Rust integration}

\paragraph{Measurements become part of parameter review.}
The benchmark plots show why the spine should include cost rather than stop at
a theorem.  Increasing salt cardinality to $k=5$ is computationally cheap in
the measured Poseidon2 configuration, though every opening carries the larger
salt.  Increasing digest width grows proof size throughout the sweep, and the
full $D=9$ binding certificate crosses a parent-hash permutation boundary.
These are ordinary engineering tradeoffs once the security target is explicit.

\paragraph{The evidence register supports ongoing development.}
The current artefact supplies the complete shared-oracle binding certificate,
the parameter constructor, typed hiding randomness, structural hiding lemmas,
and salt-cardinality capstones.  The next hiding composition follows the
Chiesa--Yogev structure: the oracle-native commitment game, basic-commitment
hiding, the bottom-up Merkle hybrid, and selective-opening simulation.
Extraction similarly has a named cache-extractor bridge.  Presenting these as
the next contracts in the spine keeps the report precise without treating an
active development sequence as a defect.

\paragraph{What this enables for \textsf{Arkworks}.}
Constructions above the commitment layer can depend on \textsf{ark-vc} without
silently choosing their own interpretation of binding, hiding, or
extractability, and without coupling themselves to one implementation.
Reviewers receive one place to inspect the property, model, parameters, Rust
obligations, measured cost, and formal evidence.  Once that layer is verified,
the remaining work above it becomes a question of composing stated contracts
rather than carrying an unexamined commitment assumption through the stack.

\takeaway{The larger result is a method for verifying \textsf{ark-vc} and
\textsf{ark-mt}.  It turns commitment security into a parameterised,
measurable, and reviewable contract, and gives later verification a stable
object to compose.}
